\documentclass[11pt]{article}
\usepackage[ngerman]{babel}
\usepackage[a4paper,margin=0.9in]{geometry}
\usepackage[T1]{fontenc}
\usepackage[utf8]{inputenc}
\usepackage{lmodern}
\usepackage{amsmath,amssymb,mathtools,mathrsfs}
\usepackage{microtype}
\usepackage{fancyhdr}
\usepackage{array,longtable,enumitem}
\setlength{\parindent}{1.5em}
\setlength{\parskip}{0.18em}
\emergencystretch=8em
\setlength{\headheight}{14pt}
\pagestyle{fancy}
\fancyhf{}
\cfoot{\thepage}
\newcommand{\sect}[2]{\section*{\S\,#1. #2}}
\newcommand{\subsect}[1]{\subsection*{#1}}
\newcommand{\perm}[2]{\begin{pmatrix}#1\\#2\end{pmatrix}}
\newcommand{\cyc}[1]{(#1)}
\newcommand{\R}{\mathfrak R}
\newcommand{\A}{\mathfrak A}
\newcommand{\N}{\mathfrak N}
\newcommand{\M}{\mathfrak M}
\newcommand{\D}{\Delta}
\newcommand{\eps}{\varepsilon}
\newcommand{\id}{\mathrm{1}}
\lhead{Weber Vol. I \S\S 177--182}
\rhead{Deutsche Quelle}
\begin{document}
\begin{center}
{\large Erster Band.}\\[0.35em]
{\Large Siebzehnter Abschnitt. Algebraische Auflösung von Gleichungen.}\\[0.35em]
{\large \S\S\,177--182}
\end{center}

\sect{177}{Einfachheit der alternirenden Gruppe}

Wir haben früher gesehen (§ 149), dass, wenn die Coefficienten einer Gleichung $n$-ten Grades als unabhängige Variable und der Körper aller rationalen Functionen dieser Coefficienten als Rationalitätsbereich betrachtet wird, die Galois'sche Gruppe der Gleichung die symmetrische Gruppe ist. In ihr ist stets ein Normaltheiler vom Index $2$ enthalten, die alternirende Gruppe, auf welche sich die Gruppe reducirt, wenn die Quadratwurzel aus der Discriminante adjungirt wird.

Bei vier Ziffern hat die alternirende Gruppe, bestehend aus der identischen Permutation, den acht dreigliedrigen Cykeln und den drei Paaren von Transpositionen, den Normaltheiler vom Index $3$
\[
  1,\\ (0,1)(2,3),\\ (0,2)(1,3),\\ (0,3)(1,2).
\]
Diese Gruppe hat ihrerseits drei Normaltheiler vom Index $2$, etwa $1,(0,1)(2,3)$; aus diesem geht durch einen weiteren Schritt die Einheitsgruppe hervor. Die Gruppe der vierundzwanzig Permutationen von vier Ziffern ist also metacyklisch, und darauf beruht die Auflösung der biquadratischen Gleichung (§§ 160, 161).

Wir weisen nun nach, dass für $n>4$ die alternirende Gruppe ausser der Einheitsgruppe überhaupt keinen Normaltheiler besitzt, also einfach ist. Daraus folgt, dass die für algebraische Auflösbarkeit nothwendige Bedingung bei Gleichungen höheren als vierten Grades, deren Gruppe symmetrisch oder alternirend ist, nicht erfüllt ist. Solche Gleichungen sind, solange die Coefficienten unabhängige Variable sind, nicht mehr algebraisch lösbar.

Sei $A$ die alternirende Gruppe der Permutationen von $n$ Ziffern $0,1,2,\ldots,n-1$ und sei $Q$ ein Normaltheiler von $A$. Die Permutationen von $A$ lassen sich aus dreigliedrigen Cykeln zusammensetzen. Ferner erhält man aus einer Permutation $\kappa$ durch Transformation mit $\pi$ die Permutation $\pi^{-1}\kappa\pi$, indem in den Cykeln von $\kappa$ die Vertauschungen $\pi$ vorgenommen werden.

Ist $\kappa$ ein dreigliedriger Cyklus, etwa $(0,1,2)$, so kann $\pi$ aus $A$ so gewählt werden, dass $\pi^{-1}\kappa\pi$ jeden beliebigen dreigliedrigen Cyklus liefert: man wählt in
\[
  \pi=\perm{0,1,2,3,\ldots,n-1}{a_0,a_1,a_2,a_3,\ldots,a_{n-1}}
\]
die drei ersten Ziffern beliebig, und vertauscht, falls es nöthig ist, noch zwei spätere Ziffern, damit $\pi$ zu $A$ gehört. Es entstehen die Cyklen $(a_0,a_1,a_2)$ oder $(a_0,a_2,a_1)$, die Potenzen voneinander sind. Enthält also $Q$ einen dreigliedrigen Cyklus, so enthält $Q$ alle, und daher $Q=A$.

Es bleibt zu zeigen, dass jeder von der Einheit verschiedene Normaltheiler $Q$ einen dreigliedrigen Cyklus enthält. Ist $\kappa\in Q$ und $\pi\in A$, so liegt auch $\pi^{-1}\kappa\pi$ in $Q$; folglich liegt auch der Commutator
\[
  \lambda=\kappa^{-1}\pi^{-1}\kappa\pi                                      \tag{1}
\]
in $Q$. Man wählt $\pi$ so, dass aus einer nicht identischen Permutation $\kappa$ ein dreigliedriger Cyklus entsteht oder eine einfachere Permutation, auf die der vorhergehende Fall anwendbar ist. Bei der Bildung von $\lambda$ dürfen diejenigen Cykeln von $\kappa$ weggelassen werden, deren Ziffern durch $\pi$ nicht geändert werden.

Die möglichen Formen von $\kappa$ sind folgende.
\begin{enumerate}[label=\arabic*.]
\item Enthält $\kappa$ einen Cyklus von mehr als drei Ziffern, etwa $(1,2,3,\ldots,m)$, so setze man $\pi=(1,2,3)$. Dann ist
\[
\lambda=(1,m,m-1,\ldots,2)(2,3,1,4,\ldots,m)=(1,2,4),
\]
also ein dreigliedriger Cyklus.
\item Enthält $\kappa$ zwei dreigliedrige Cykeln,
\[
  \kappa=(1,2,3)(4,5,6),
\]
so setze man $\pi=(1,3,4)$ und erhält
\[
\lambda=(1,3,2)(4,6,5)(3,2,4)(1,5,6)=(1,2,5,3,4),
\]
was unter Fall 1 fällt.
\item Enthält $\kappa$ einen dreigliedrigen und einen zweigliedrigen Cyklus,
\[
  \kappa=(1,2,3)(4,5),
\]
so ergiebt $\pi=(1,2,4)$
\[
\lambda=(1,3,2)(4,5)(2,4,3)(1,5)=(1,2,5,3,4),
\]
wieder Fall 1.
\item Enthält $\kappa$ drei Transpositionen,
\[
  \kappa=(1,2)(3,4)(5,6),
\]
so folgt aus $\pi=(1,3,5)$
\[
\lambda=(1,2)(3,4)(5,6)(3,2)(5,4)(1,6)=(1,3,5)(2,6,4),
\]
also Fall 2.
\item Enthält $\kappa$ zwei Transpositionen und ein unverändertes Element,
\[
  \kappa=(1,2)(3,4)(5),
\]
so setzt man $\pi=(1,2,5)$ und erhält
\[
\lambda=(1,2)(3,4)(5)(2,5)(3,4)(1)=(1,5,2).
\]
\end{enumerate}
Damit sind für $n>4$ alle Fälle erschöpft. Für $n=4$ bleibt der Fall zweier Transpositionen allein übrig; gerade dieser Ausnahmefall ermöglicht die algebraische Auflösung der Gleichung vierten Grades.

Aus dem Satze folgt weiter, dass die symmetrische Gruppe keine anderen Normaltheiler besitzt als sich selbst, die alternirende Gruppe und die Einheitsgruppe. Ist nämlich $S$ symmetrisch, $A$ alternirend und $Q$ ein Normaltheiler von $S$, so ist $Q'=A\cap Q$ ein Normaltheiler von $A$, also $A$ oder $1$. Im ersten Falle ist $Q=A$ oder $S$. Im zweiten Falle enthält $Q$ ausser der Einheit keine Permutation erster Art; zwei verschiedene nicht identische Elemente von $Q$ können dann nicht vorkommen, und ein einzelnes solches Element müsste durch jede Umbenennung der Ziffern unverändert bleiben. Das ist nur bei zwei Ziffern möglich, wo $1,(1,2)$ schon die ganze symmetrische Gruppe ist.

\sect{178}{Nicht metacyklische Gleichungen im Körper der rationalen Zahlen}

Durch den Satz des vorigen Paragraphen ist gezeigt, dass eine Gleichung $n$-ten Grades für $n>4$ nicht mehr algebraisch gelöst werden kann, wenn ihre Coefficienten unabhängige Veränderliche sind. Von noch grösserem Interesse ist die Frage, ob es im Körper der rationalen Zahlen Gleichungen $n$-ten Grades giebt, die nicht algebraisch lösbar sind; allgemeiner, ob ganzzahlige Gleichungen vorkommen, deren Gruppe die symmetrische Gruppe ist, die also keinen Affect haben.

Bildet man die Galois'sche Resolvente $G(t)=0$ vom Grade $\nu(n)$ einer allgemeinen Gleichung $n$-ten Grades $f(x)=0$ mit unbestimmten Coefficienten $a$, so ist $G(t)$ eine ganze Function der Veränderlichen $t$ und $a$, welche nicht in Factoren rationaler Functionen von $t$ und $a$ zerfällt. Werden für die Variablen $a$ Grössen eines Körpers $\Omega$ eingesetzt und wird $G(t)$ in diesem Körper reducibel, so hat $f(x)=0$ im Rationalitätsbereich $\Omega$ einen Affect. Die Frage ist also, ob man die $a$ rational so wählen kann, dass $G(t)$ über den rationalen Zahlen irreducibel bleibt.

Hilbert hat hierzu den allgemeinen Satz bewiesen, dass man in einer irreduciblen Function beliebig vieler Variablen für einen beliebigen Theil der Variablen rationale Zahlen so einsetzen kann, dass eine irreducible Function der übrigen Variablen entsteht. Weber benutzt hier nur eine einfachere Folgerung: für jeden Primzahlgrad lassen sich Gleichungen ohne Affect finden.

Wir haben in § 153, 9 bewiesen, dass eine transitive Permutationsgruppe von $n$ Ziffern, die nicht die symmetrische Gruppe ist, bei primem $n$ keine einzelne Transposition enthalten kann. Hat also die irreducible Gleichung $f(x)=0$ von Primzahlgrad $n$ einen Affect, so ist ihre Gruppe transitiv, aber ohne Transposition. Werden beliebige $n-2$ Wurzeln adjungirt, so könnte ausser der Einheitspermutation nur die Vertauschung der beiden letzten Wurzeln übrig bleiben; diese ist aber eine Transposition und kommt nicht vor. Daher:

\begin{quote}
1. Wenn eine irreducible Gleichung, deren Grad $n$ eine Primzahl ist, einen Affect hat, so können zwei beliebige ihrer Wurzeln rational durch die übrigen ausgedrückt werden.
\end{quote}

Daraus folgt:
\begin{quote}
2. Wenn der Körper $\Omega$ nur reelle Zahlen enthält, so kann eine irreducible Gleichung von Primzahlgrad $n$ mit einem Affect in $\Omega$ nicht zwei imaginäre und $n-2$ reelle Wurzeln haben.
\end{quote}

Es giebt aber Gleichungen beliebigen Grades mit reellen Coefficienten, die genau zwei conjugirt imaginäre und $n-2$ reelle Wurzeln haben. In
\[
  f(x)=(x-\alpha_1)(x-\alpha_2)\cdots(x-\alpha_n)
\]
wähle man $\alpha_1,\alpha_2$ conjugirt imaginär und die übrigen $\alpha$ reell. Da die Anzahl reeller und imaginärer Wurzeln bei hinreichend kleinen stetigen Aenderungen der Coefficienten unverändert bleibt, kann man die Coefficienten rational beliebig nahe wählen. Es bleibt, die Irreducibilität zu sichern.

Dazu dient das Eisenstein'sche Kriterium in folgender Form. Ist $p$ eine Primzahl und sind $c_0,c_1,\ldots,c_n$ ganze Zahlen, wobei $c_0$ nicht durch $p$, die $c_1,\\ldots,c_n$ durch $p$, aber $c_n$ nicht durch $p^2$ theilbar ist, so ist
\[
  \varphi(x)=c_0x^n+c_1x^{n-1}+\cdots+c_n
\]
irreducibel. Denn zerfiele $\varphi$ in zwei Factoren mit rationalen, also nach § 2 auch ganzzahligen Coefficienten,
\[
\varphi(x)=(a_0x^h+a_1x^{h-1}+\cdots+a_h)
          (\beta_0x^k+\beta_1x^{k-1}+\cdots+\beta_k),
\]
mit $h,k>0$, so wäre $a_h\beta_k=c_n$. Einer der beiden Factoren $a_h,\beta_k$ wäre durch $p$ theilbar, der andere nicht. Sei etwa $\beta_k$ durch $p$ theilbar, $a_h$ nicht, und sei $\beta_v$ der letzte nicht durch $p$ theilbare Coefficient unter den $\beta$. Dann ist der Coefficient von $x^{k-v}$ im Product nicht durch $p$ theilbar, während er nach der Voraussetzung durch $p$ theilbar sein müsste. Widerspruch.

Setzt man für die Coefficienten von $f(x)$
\[
  a_1={c_1\over c_0},\quad a_2={c_2\over c_0},\quad \ldots,\quad a_n={c_n\over c_0},
\]
so bleibt in der Wahl der ganzen Zahlen $c_i$ genug Freiheit, um die rationalen Brüche $a_i$ einem gegebenen Werthsystem beliebig nahe zu bringen und gleichzeitig Irreducibilität zu erzwingen. Damit ist bewiesen:
\begin{quote}
A. Es giebt von jedem beliebigen Primzahlgrade unendlich viele Gleichungen mit rationalen Coefficienten ohne Affect.
\end{quote}

\sect{179}{Auflösung durch reelle Radicale}

Bei der cubischen Gleichung mit reellen Coefficienten unterscheidet man seit Alters die Fälle negativer und positiver Discriminante. Im zweiten Falle hat die Gleichung drei reelle Wurzeln, während die Cardanische Formel die dritte Wurzel aus einer imaginären Grösse verlangt; der Versuch, eine rein reelle Darstellung zu erzwingen, führt wieder auf eine cubische Gleichung derselben Art. Dieser Fall heisst daher der \emph{casus irreducibilis}. Er ist ein besonderer Fall der cyklischen Gleichungen mit reellen Wurzeln, bei denen Wurzeln aus imaginären Grössen, oder äquivalent Winkeltheilungen, auftreten. Dass diese imaginären Grössen nicht vermieden werden können, wird nun bewiesen.

Sei $\Omega$ ein reeller Rationalitätsbereich und $g(t)=0$ eine Normalgleichung in $\Omega$. Ist eine Wurzel reell, so sind alle Wurzeln reell; andernfalls treten die Wurzeln paarweise conjugirt imaginär auf. Ist $g(t)$ die Galois'sche Resolvente einer Gleichung $F(x)=0$, so sind die Wurzeln von $F$ rational durch eine Wurzel von $g$ darstellbar, und daher hat $g$ reelle Wurzeln genau dann, wenn auch $F$ nur reelle Wurzeln hat.

Wird $g(t)$ durch Adjunction einer Wurzel $\theta$ einer irreduciblen Gleichung $\chi=0$ von Primzahlgrad reducirt, so liegen nach § 157 alle Wurzeln von $\chi$ in $\Omega(\theta)$. Ist also $\theta$ reell, so sind alle Wurzeln von $\chi$ reell. Folglich:
\begin{quote}
1. Eine Normalgleichung mit reellen Wurzeln kann nur durch solche irreducible Gleichungen von Primzahlgrad reducirt werden, die lauter reelle Wurzeln haben.
\end{quote}

Wir fragen nun, ob eine Reduction der Normalgleichung durch Adjunction eines reellen Radicals $\sqrt[p]{a}$ bewirkt werden kann. Man darf $p$ als Primzahl voraussetzen. Ist $a$ keine $p$-te Potenz in $\Omega$, so ist
\[
  x^p-a=0                                                        \tag{2}
\]
irreducibel. Denn sind $\alpha=\sqrt[p]{a}$ und $\varepsilon$ eine primitive $p$-te Einheitswurzel, so sind die Wurzeln
\[
  \alpha,\ \varepsilon\alpha,\ \varepsilon^2\alpha,\ldots,\varepsilon^{p-1}\alpha. \tag{1}
\]
Zerfiele
\[
  x^p-a=f_1(x)f_2(x),                                           \tag{3}
\]
so müsste für einen Factor $f_1$ vom Grade $\nu$, $0<\nu<p$, das Product einer Auswahl von $\nu$ dieser Wurzeln eine Grösse $b$ aus $\Omega$ sein; also
\[
  a^\nu=b^p.                                                    \tag{4}
\]
Da $(\nu,p)=1$, gibt es ganze Zahlen $h,k$ mit $\nu h+pk=1$. Dann wäre
\[
  a=a^{\nu h+pk}=(b^h a^k)^p,
\]
gegen die Voraussetzung.

Nehmen wir nun $\Omega$ reell an. Für ungerades $p$ hat (2) imaginäre Wurzeln und kann daher nach 1. eine Normalgleichung mit reellen Wurzeln nicht reducieren. Für $p=2$ kann dies nur geschehen, wenn $a>0$ und der Grad von $g(t)$ gerade ist. Daher:
\begin{quote}
2. Eine Normalgleichung ungeraden Grades kann nicht durch Adjunction eines reellen Radicals reducirt werden.
\end{quote}

Dies trifft den casus irreducibilis der cubischen Gleichung nach Adjunction der Quadratwurzel der positiven Discriminante. Der Rationalitätsbereich bleibt reell und bleibt es nach Adjunction reeller Radicale; die Normalgleichung hat ungeraden Grad und kann daher nicht reducirt werden.

Dementsprechend zerfallen die irreduciblen cubischen Gleichungen über den rationalen Zahlen in die cyklischen Gleichungen vom Grade $3$ mit drei reellen Wurzeln, die nicht durch reelle Radicale lösbar sind, sowie in Gleichungen mit Gruppe sechsten Grades und positiver oder negativer Discriminante. Die positive Discriminante führt zum casus irreducibilis und zur Dreitheilung eines Winkels; die negative Discriminante umfasst als besondere Fälle die reinen Gleichungen $x^3=a$, etwa $x^3=2$ für die Würfelverdoppelung.

\sect{180}{Metacyklische Gleichungen von Primzahlgrad}

Die allgemeinen Bedingungen des § 176 sind für die unmittelbare Anwendung noch zu unhandlich. Wir leiten daher, zunächst für Gleichungen von Primzahlgrad $n>2$, das von Galois aufgestellte Kriterium her.

Sei $f(x)=0$ irreducibel vom Primzahlgrad $n$ und sei $P$ ihre Gruppe. Soll $f$ algebraisch lösbar sein, so muss $P$ metacyklisch sein; es gibt also eine Kette
\[
  P,\ P_1,\ P_2,\ldots, P_{\mu-1},\ 1,                         \tag{1}
\]
in der jedes Glied Normaltheiler des vorhergehenden von Primzahlindex ist. Die Gleichung selbst kann vor der letzten Adjunction nicht reducirt werden; nach der letzten zerfällt sie in $n$ lineare Factoren. Daher ist $P_{\mu-1}$ eine cyklische Gruppe vom Grade $n$. Durch geeignete Numerirung der Wurzeln kann ihr Erzeuger als
\[
  \pi=(0,1,2,\ldots,n-1)
\]
geschrieben werden. Bezeichnet man die Wurzeln mit $x_z$ und setzt $x_{z+n}=x_z$, so besteht $P_{\mu-1}$ aus den Substitutionen
\[
  (z,z+h),\qquad h=0,1,2,\ldots,n-1.                            \tag{2}
\]

Diese sind besondere Fälle der linearen Substitution
\[
  (z,az+b),                                                     \tag{3}
\]
wobei $a\not\equiv0\pmod n$ ist. Die Gesamtheit aller solcher Permutationen bildet eine Gruppe vom Grade $n(n-1)$, die lineare Gruppe. Sind
\[
  \lambda=(z,az+b),\qquad \lambda'=(z,a'z+b'),
\]
so gilt
\[
  \lambda\lambda'=(z,aa'z+a'b+b').                              \tag{4}
\]
Ist $g$ eine primitive Wurzel modulo $n$ und $a_0$ ein fester Vertreter, so besteht jede transitive lineare Gruppe aus den Substitutionen
\[
  \lambda=(z,a_0^h z+b),\qquad h=0,1,\ldots,e-1,
  \quad b=0,1,\ldots,n-1,                                      \tag{5}
\]
wobei $e\mid n-1$. Umgekehrt bildet jedes System dieser Form eine Gruppe.

Der entscheidende Satz lautet:
\begin{quote}
I. Ist eine transitive lineare Gruppe $L$ Normaltheiler einer anderen Permutationsgruppe $P$ derselben $n$ Ziffern, so ist auch $P$ eine lineare Gruppe.
\end{quote}

Zum Beweis stellt man eine beliebige Permutation $\pi$ durch eine Function $\varphi(z)$ dar:
\[
 \pi=\perm{0,1,2,\ldots,n-1}{a_0,a_1,a_2,\ldots,a_{n-1}},
 \qquad a_z=\varphi(z).
\]
Nach Lagrange'scher Interpolation, modulo $n$, kann man schreiben
\[
  \varphi(z)\equiv -a_0 z^{n-1}-a_1(z-1)^{n-1}-\cdots-a_{n-1}(z-n+1)^{n-1}\pmod n. \tag{6}
\]
Ist $L\triangleleft P$ und enthält $L$ die cyklische Substitution $(z,z+1)$, so muss für eine Permutation $\pi=(z,\varphi(z))$ aus $P$ eine lineare Substitution $(z,a'z+a)$ existieren, so dass
\[
  \varphi(z+1)\equiv a'\varphi(z)+a\pmod n.                     \tag{7}
\]
Da $\varphi$ den Grad $n$ nicht erreicht, folgt durch Vergleich der höchsten Potenzen $a'\equiv1$, also
\[
  \varphi(z+h)\equiv \varphi(z)+ah,
\]
und daher
\[
  \varphi(z)=az+b.                                               \tag{8}
\]
Somit ist jede Permutation von $P$ linear.

Aus diesem Satze folgt der Galois'sche Satz:
\begin{quote}
II. Die Gruppe einer metacyklischen Gleichung von Primzahlgrad ist linear.
\end{quote}
Denn $P_{\mu-1}$ ist cyklisch, also linear; da es Normaltheiler von $P_{\mu-2}$ ist, ist $P_{\mu-2}$ linear, und so aufwärts bis $P$.

Umgekehrt gilt:
\begin{quote}
III. Jede irreducible Gleichung von Primzahlgrad, deren Gruppe linear ist, ist metacyklisch.
\end{quote}
Denn jede transitive lineare Gruppe $L$ besitzt, falls sie nicht schon cyklisch ist, einen Normaltheiler $L'$ von Primzahlindex, der wieder transitiv linear ist. In der Darstellung (5) erhält man für einen Primteiler $p$ von $e=pe'$ die Gruppe
\[
  L'=\bigl\{(z,a_0^{ph}z+b):\ h=0,1,\ldots,e'-1,
       \ b=0,1,\ldots,n-1\bigr\}.
\]

Die volle lineare Gruppe wird durch die Substitutionen
\[
  s=(z,z+1),\qquad t=(z,gz)                                      \tag{9}
\]
erzeugt; ein Theiler der Form (5) durch
\[
  s=(z,z+1),\qquad t_0=(z,a_0z).                                \tag{10}
\]
Nimmt man $a_0=g^2$, so erhält man die halbmetacyklische Gruppe Kroneckers. Die volle metacyklische Gruppe ist kein Theiler der alternirenden Gruppe; der grösste gemeinschaftliche Theiler der beiden ist die halbmetacyklische Gruppe.

Man kann die Bedingung der Lösbarkeit auch in Resolventenform fassen. Die volle lineare Gruppe ist ein Theiler vom Index
\[
  \nu=1\cdot2\cdot3\cdots(n-2)
\]
in der symmetrischen Gruppe. Eine zu ihr gehörige Function $y$ genügt einer Resolvente $F(y)=0$ vom Grade $\nu$. Ist eine Wurzel dieser Resolvente rational und keine Doppelwurzel, so liegt die Gruppe der Gleichung in einer conjugirten metacyklischen Gruppe. Daher:
\begin{quote}
V. Die nothwendige und hinreichende Bedingung für die Lösbarkeit von $f(x)=0$ durch Radicale ist, dass die Resolvente $\nu$-ten Grades $F(y)=0$, passend ohne Doppelwurzeln gewählt, eine rationale Wurzel besitzt.
\end{quote}

Eine andere Form lautet:
\begin{quote}
VI. Ist eine irreducible Gleichung von Primzahlgrad metacyklisch, so sind alle Wurzeln rational durch zwei beliebige unter ihnen ausdrückbar.
\end{quote}
Denn in einer linearen Gruppe lässt ausser der Identität keine Permutation zwei Ziffern fest. Der Satz gilt auch umgekehrt. Hat eine irreducible Gleichung vom Primzahlgrad die Eigenschaft, dass alle Wurzeln rational durch zwei beliebige ausgedrückt werden können, so kann keine nicht identische Permutation der Gruppe zwei Ziffern festlassen. Zählt man die Permutationen, welche eine bestimmte Ziffer festlassen, und die cyklischen Permutationen $n$-ten Grades, so erhält man eine cyklische Gruppe $C=1,\gamma,\ldots,\gamma^{n-1}$ als Normaltheiler; nach Satz I ist die ganze Gruppe linear. Also:
\begin{quote}
VII. Eine irreducible Gleichung von Primzahlgrad, bei der alle Wurzeln rational durch zwei beliebige von ihnen ausdrückbar sind, ist metacyklisch.
\end{quote}

Für reelle Coefficienten folgt weiter:
\begin{quote}
VIII. Eine irreducible metacyklische Gleichung von ungeradem Primzahlgrad mit reellen Coefficienten hat entweder lauter reelle Wurzeln oder nur eine.
\end{quote}
Denn zwei reelle Wurzeln würden nach VI alle Wurzeln reell machen. Das Vorzeichen der Discriminante ist im zweiten Falle $(-1)^{(n-1)/2}$. Daher:
\begin{quote}
IX. Wenn $n\equiv1\pmod4$ ist, so ist die Discriminante immer positiv; ist aber $n\equiv3\pmod4$, so entscheidet das Vorzeichen der Discriminante, welcher der beiden Fälle des Satzes VIII eintritt.
\end{quote}

\sect{181}{Anwendung auf die metacyklischen Gleichungen fünften Grades}

Für $n=5$ hat die cyklische Gruppe $5$, die halbmetacyklische $10$ und die volle metacyklische Gruppe $20$ Permutationen. Eine metacyklische Function genügt bei der allgemeinen Gleichung fünften Grades einer Resolvente sechsten Grades; hat diese eine rationale Wurzel, so ist die Gleichung fünften Grades metacyklisch. Die halbmetacyklischen Functionen genügen einer Resolvente zwölften Grades, die nach Adjunction von $\sqrt\Delta$ in zwei Factoren sechsten Grades zerfällt.

Für den Modul $5$ lauten die erzeugenden Substitutionen
\[
  s=(z,z+1),\qquad t=(z,2z),\qquad t^2=(z,4z).
\]
Auf die Ziffern $0,1,2,3,4$ wirken sie durch
\[
\begin{array}{c|ccccc}
     &0&1&2&3&4\\ \hline
s    &1&2&3&4&0\\
t    &0&2&4&1&3\\
t^2  &0&4&3&2&1
\end{array}                                                       \tag{1}
\]
und auf die fünf Paare $(0,1),(1,2),(2,3),(3,4),(4,0)$ durch
\[
\begin{array}{c|ccccc}
      &(0,1)&(1,2)&(2,3)&(3,4)&(4,0)\\ \hline
s     &(1,2)&(2,3)&(3,4)&(4,0)&(0,1)\\
t^2   &(0,4)&(4,3)&(3,2)&(2,1)&(1,0)
\end{array}                                                       \tag{2}
\]
und
\[
\begin{array}{c|ccccc}
      &(0,1)&(1,2)&(2,3)&(3,4)&(4,0)\\ \hline
t     &(0,2)&(2,4)&(4,1)&(1,3)&(3,0).
\end{array}                                                       \tag{3}
\]
Man kann halbmetacyklische Functionen auf verschiedene Arten bilden. Die einfachste ist
\[
  u=x_0x_1+x_1x_2+x_2x_3+x_3x_4+x_4x_0,                         \tag{4}
\]
die durch $t$ in
\[
  u'=x_0x_2+x_2x_4+x_4x_1+x_1x_3+x_3x_0                          \tag{5}
\]
übergeht. Ist
\[
  f(x)=x^5-ax^4+bx^3-cx^2+dx-e=0,                                \tag{6}
\]
so ist
\[
  u+u'=b,                                                         \tag{7}
\]
und
\[
  y=u-u'                                                          \tag{8}
\]
ist halbmetacyklisch, während $y^2$ vollmetacyklisch ist. Ist $\sqrt\Delta$ das Product der zehn Wurzeldifferenzen, so ist
\[
  Y={u-u'\over\sqrt\Delta}                                       \tag{9}
\]
ebenfalls vollmetacyklisch und Wurzel einer Gleichung sechsten Grades.

Bezeichnet man die vollmetacyklische Gruppe mit $M$, die halbmetacyklische mit $N$ und die alternirende mit $A$, so zerlegt sich
\[
 A=N+N(1,2)(3,4)+Nt(0,1)+Nt(0,2)+Nt(0,3)+Nt(0,4).                \tag{10}
\]
Hieraus entstehen für $u$ sechs innerhalb $A$ conjugirte Werthe $u_1,\ldots,u_6$; die entsprechenden $u'_1,\ldots,u'_6$ erhält man durch Anwendung von $t$, oder dadurch, dass man in jedem $u_i$ die fehlenden Paare vereinigt. Für alle gilt $u_i+u'_i=b$.

Die sechs Grössen $y_i=u_i-u'_i$ sind die Wurzeln einer Gleichung sechsten Grades, deren Coefficienten rational von den Coefficienten von $f$ und von $\sqrt\Delta$ abhängen. Da alle $y_i$ bei Wechsel des Vorzeichens von $\sqrt\Delta$ ihr Vorzeichen ändern, hat sie die Form
\[
 y^6+a_2y^4+a_4y^2+a_6-
 \sqrt\Delta\,(a_1y^5+a_3y^3+a_5y)=0.                            \tag{12}
\]
Die Coefficienten $a_i$ sind ganze Functionen der Coefficienten von $f$. Durch Gradbetrachtung folgt $a_1=a_3=0$, und $a_5$ ist eine Zahl.

Weber bildet die Resolvente für die Bring-Jerrard'sche Form
\[
  x^5+\alpha x+\beta=0.                                          \tag{13}
\]
Dann hängen $a_2,a_4,a_6$ nicht von $\beta$ ab, und man hat
\[
  a_2=m_1\alpha,
  \qquad a_4=m_2\alpha^2,
  \qquad a_6=m_3\alpha^3,
  \qquad a_5=m.                                                   \tag{14}
\]
Setzt man $\beta=0$, so wird
\[
 x_0=0,\\ x_1=\sqrt{-\alpha},\\ x_2=i\sqrt{-\alpha},
 \qquad x_3=-\sqrt{-\alpha},\\ x_4=-i\sqrt{-\alpha},              \tag{15}
\]
und
\[
  \sqrt\Delta=-16\sqrt{\alpha^5},
  \qquad \Delta=256\alpha^5.                                     \tag{16}
\]
Hieraus erhält man für $\beta=0$
\[
\begin{aligned}
&y^6+a_2y^4+a_4y^2+a_6-a_5\sqrt\Delta\,y\\
&=(y+2\sqrt\alpha)^4(y^2-8y\sqrt\alpha+20\alpha)\\
&=y^6-20\alpha y^4+240\alpha^2 y^2
  +512\sqrt{\alpha^5}\,y+320\alpha^3 .
\end{aligned}                                                     \tag{17}
\]
Daraus folgt allgemein
\[
  y^6-20\alpha y^4+240\alpha^2y^2
  -32\sqrt\Delta\,y+320\alpha^3=0.                              \tag{18}
\]
Die Discriminante ist
\[
  \Delta=5^5\beta^4+2^8\alpha^5.                                \tag{19}
\]
Für $u$ wird die Gleichung einfacher:
\[
  u^6-5\alpha u^4+15\alpha^2u^2-
  \sqrt\Delta\,u+5\alpha^3=0.                                    \tag{20}
\]
Mit $u^2=v$ ergibt sich die rationale Gleichung sechsten Grades
\[
  \bigl(v^3-5\alpha v^2+15\alpha^2v+5\alpha^3\bigr)^2
  =\Delta v,                                                     \tag{21}
\]
oder, nach Umformung,
\[
  (v-
\alpha)^4(v^2+6\alpha v+25\alpha^2)=5^5\beta^4 v.           \tag{22}
\]

Für eine irreducible Gleichung der Form (13) genügt also zur Prüfung der Auflösbarkeit, die Irreducibilität festzustellen und dann zu untersuchen, ob (21) oder (22) eine rationale Wurzel hat. Sind $\alpha,\beta$ ganze Zahlen, so ist auch ein rationales $v$ ganzzahlig und unter den Factoren von $25\alpha^6$ zu suchen. Insbesondere ist keine Gleichung
\[
  x^5+5x+5t=0
\]
metacyklisch, wenn $t$ nicht durch $5$ theilbar ist.

Um metacyklische Gleichungen zu ermitteln, setzt man in (22)
\[
  \beta=\alpha\mu,
  \qquad v=\alpha\lambda,
\]
und erhält
\[
  \alpha={5^5\mu^4\lambda\over(\lambda-1)^4(\lambda^2+6\lambda+25)},
  \qquad
  \beta={5^5\mu^5\lambda\over(\lambda-1)^4(\lambda^2+6\lambda+25)}. \tag{23}
\]
Werden hierin für $\lambda$ und $\mu$ rationale Grössen gesetzt, so ist
\[
  x^5+\alpha x+
\beta=0                                                     \tag{24}
\]
algebraisch lösbar. Umgekehrt ist keine irreducible Gleichung der Form (24) algebraisch lösbar, wenn $\alpha,
\beta$ nicht in der Form (23) darstellbar sind. Zum Beispiel liefern $\lambda=-1$, $\mu=1$ nach der Substitution $x\xi=5$ die irreducible lösbare Gleichung
\[
  \xi^5+5\xi^4-5\cdot64=0.
\]

\sect{182}{Die Gruppe der Resolvente}

Die sechs Grössen des vorigen Paragraphen
\[
\begin{array}{lll}
 v_1=(u_1-u'_1)^2,& v_2=(u_2-u'_2)^2,& v_3=(u_3-u'_3)^2,\\
 v_4=(u_4-u'_4)^2,& v_5=(u_5-u'_5)^2,& v_6=(u_6-u'_6)^2
\end{array}                                                       \tag{1}
\]
sind die Wurzeln einer Gleichung sechsten Grades $F(v)=0$, deren Coefficienten rational durch die symmetrischen Grundfunctionen $a,b,c,d,e$ der $x$ ausdrückbar sind. Diese Gleichung ist eine Totalresolvente der Gleichung fünften Grades. Ihre Gruppe hat denselben Grad wie die Gruppe von $f(x)=0$, also $120$. Da die symmetrische Gruppe von sechs Ziffern den Grad $720$ hat, folgt:
\begin{quote}
Die symmetrische Permutationsgruppe von sechs Ziffern hat einen transitiven Divisor vom Index $6$.
\end{quote}

Um diese Gruppe $C$ zu beschreiben, untersucht Weber den Einfluss der vier Transpositionen $(0,1)$, $(0,2)$, $(0,3)$, $(0,4)$ der fünf Wurzeln auf die Indices der $v$. Man erhält die erzeugenden Permutationen
\[
\begin{array}{rcl}
(0,1):&\pi_1&=(1,3)(2,5)(4,6),\\
(0,2):&\pi_2&=(1,4)(2,3)(5,6),\\
(0,3):&\pi_3&=(1,5)(2,6)(3,4),\\
(0,4):&\pi_4&=(1,6)(2,4)(3,5).
\end{array}
\]
Unter den $120$ Permutationen der Gruppe $C$ kommen zum Beispiel vor
\[
\begin{aligned}
\pi_1\pi_2&=(1,2,6)(3,4,5),\\
\pi_1\pi_2\pi_3&=(1,6,5,4),\\
\pi_1\pi_2\pi_3\pi_4&=(2,4,6,3,5),\\
(\pi_1\pi_2\pi_3\pi_4)^3&=(2,3,4,5,6).
\end{aligned}
\]

Eine Function
\[
  v_1v_3+v_2v_4+v_5v_6
\]
bleibt durch $\pi_1$ und durch $\pi_1\pi_2$ ungeändert. Wendet man die Potenzen des Cyklus $(2,3,4,5,6)$ an, so erhält man die fünf Functionen
\[
\begin{array}{rcl}
 w_0&=&v_1v_2+v_4v_5+v_3v_6,\\
 w_1&=&v_1v_4+v_2v_6+v_3v_5,\\
 w_2&=&v_1v_6+v_2v_5+v_3v_4,\\
 w_3&=&v_1v_3+v_2v_4+v_5v_6,\\
 w_4&=&v_1v_5+v_2v_3+v_4v_6.
\end{array}                                                       \tag{2}
\]
Die Permutationen $\pi_1,
\pi_2,
\pi_3,
\pi_4$ vertauschen diese Functionen wie die Transpositionen $(0,1)$, $(0,2)$, $(0,3)$, $(0,4)$ der Indices $0,1,2,3,4$. Der Gruppe $C$ entspricht daher die symmetrische Gruppe der Indices der $w$. Eine symmetrische Function der fünf $w$, die nicht zugleich symmetrische Function der sechs $v$ ist, gehört zu $C$. Man kann etwa
\[
  W=w_0^2+w_1^2+w_2^2+w_3^2+w_4^2
\]
oder
\[
  (\lambda-w_0)(\lambda-w_1)(\lambda-w_2)(\lambda-w_3)(\lambda-w_4)
\]
nehmen. Eine solche Grösse $W$ ist die Wurzel einer irreduciblen Gleichung sechsten Grades, deren Coefficienten symmetrische Functionen der sechs Grössen $v$ sind.

\end{document}
